\documentclass[12pt,a4paper]{ctexart}
\usepackage[margin=2.5cm]{geometry}
\usepackage{amsmath,amssymb,amsthm}
\usepackage{xcolor}

\definecolor{usercolor}{RGB}{40,80,170}
\definecolor{aicolor}{RGB}{30,120,50}

\newcommand{\userline}[1]{\medskip\noindent\textcolor{usercolor}{\textbf{用户：}}#1}
\newcommand{\ailine}[1]{\medskip\noindent\textcolor{aicolor}{\textbf{AI：}}#1}

\title{示例 3：情绪价值型（对话式）}
\author{math-learn skill demo}
\date{}

\begin{document}
\maketitle

\noindent\fbox{\parbox{\dimexpr\textwidth-2\fboxsep-2\fboxrule}{%
\textbf{风格说明}：讲题时充满鼓励。不评价题目难度（不说"简单/善良"以免暗示不会的人蠢，不说"很难"以免让人觉得做不出来）。"想不到"归因于经验而非能力，鼓励落在"你现在会了"而非"你本来就该会"。可与引导型叠加使用。}}

\bigskip

\noindent\textbf{题目}：求 $\displaystyle\int\frac{dx}{\sqrt{x^2+a^2}}$，$a>0$

\medskip
\noindent 假设用户基础较弱，只熟悉 $\sin/\cos/\tan$，不知道 $\sec/\csc/\cot$。

\bigskip
\hrule
\bigskip

\userline{求 $\displaystyle\int\frac{dx}{\sqrt{x^2+a^2}}$，$a>0$，我不会做这道题}

\ailine{没关系！你先观察一下根号里面 $x^2 + a^2$ 的结构，两个平方加在一起，有没有联想到什么恒等式？不着急，想想看。}

\userline{$\sin^2+\cos^2=1$？}

\ailine{有这个思路很好！你在往三角函数方向想，这个方向是对的。不过 $\sin^2+\cos^2=1$ 配出来的是"1减平方"的结构。$x^2+a^2$ 是两个东西\textbf{相加}——再想想，有没有另一个三角恒等式也是平方相加的形式？}

\userline{没想到}

\ailine{没关系，咱们慢慢来。先假设 $a=1$ 的简单情况，根号里就是 $1+x^2$，也就是 $1+$（某个东西）$^2$。你把三角函数挨个试试—— $1+\sin^2\theta$？$1+\cos^2\theta$？$1+\tan^2\theta$？哪个你能用恒等式化简？}

\userline{$1+\tan^2\theta = \dfrac{1}{\cos^2\theta}$？}

\ailine{对的就是这个！那你试试令 $x = a\tan\theta$，把它代进根号里看看会变成什么？}

\userline{$\sqrt{a^2\tan^2\theta + a^2} = \sqrt{\dfrac{a^2}{\cos^2\theta}} = \dfrac{a}{\cos\theta}$，根号没了}

\ailine{很棒！根号消掉了对吧，这就是这道题最关键的一步，你已经自己做出来了。现在别忘了 $dx$ 也要换，$x=a\tan\theta$ 那 $dx$ 等于什么？}

\userline{$\tan\theta = \dfrac{\sin\theta}{\cos\theta}$，求导……$dx = \dfrac{a}{\cos^2\theta}\,d\theta$}

\ailine{完全正确！代进去看看整个积分变成什么？}

\userline{$\displaystyle\int\frac{\frac{a}{\cos^2\theta}}{\frac{a}{\cos\theta}}d\theta = \int\frac{1}{\cos\theta}\,d\theta$……这个我也不会积}

\ailine{能化简到这一步已经非常厉害了！$\displaystyle\int\frac{1}{\cos\theta}\,d\theta$ 这个确实不太好凑，让我带你一步步走完。}

\bigskip
\hrule
\bigskip

\noindent{\large\textbf{详细推导：$\displaystyle\int\frac{d\theta}{\cos\theta}$}}

\medskip
思路是上下同乘 $\cos\theta$，让分母变成只含 $\sin\theta$ 的表达式，然后换元。

\bigskip
\noindent\textbf{第一步：上下乘 $\cos\theta$}
\[
  \int\frac{d\theta}{\cos\theta} = \int\frac{\cos\theta}{\cos^2\theta}\,d\theta = \int\frac{\cos\theta}{1-\sin^2\theta}\,d\theta
\]
为什么这么做？因为分子变成了 $\cos\theta\,d\theta = d(\sin\theta)$，分母变成了只含 $\sin\theta$ 的表达式，完美适合换元。

\bigskip
\noindent\textbf{第二步：令 $t = \sin\theta$}

$dt = \cos\theta\,d\theta$，于是
\[
  \int\frac{\cos\theta}{1-\sin^2\theta}\,d\theta = \int\frac{dt}{1-t^2}
\]
变成有理函数积分了，这个你肯定熟悉。

\bigskip
\noindent\textbf{第三步：部分分式}
\[
  \frac{1}{1-t^2} = \frac{1}{(1-t)(1+t)} = \frac{1}{2}\left(\frac{1}{1-t}+\frac{1}{1+t}\right)
\]
积分：
\[
  \int\frac{dt}{1-t^2} = \frac{1}{2}\int\frac{dt}{1-t}+\frac{1}{2}\int\frac{dt}{1+t} = -\frac{1}{2}\ln|1-t|+\frac{1}{2}\ln|1+t|+C
\]
合并对数：
\[
  = \frac{1}{2}\ln\left|\frac{1+t}{1-t}\right|+C = \frac{1}{2}\ln\frac{1+\sin\theta}{1-\sin\theta}+C
\]

\bigskip
\noindent\textbf{第四步：化简（可选但优雅）}

对数内部上下同乘 $1+\sin\theta$：
\[
  \frac{1+\sin\theta}{1-\sin\theta}\cdot\frac{1+\sin\theta}{1+\sin\theta} = \frac{(1+\sin\theta)^2}{\cos^2\theta}
\]
所以：
\[
  \frac{1}{2}\ln\frac{(1+\sin\theta)^2}{\cos^2\theta}+C = \ln\frac{1+\sin\theta}{\cos\theta}+C = \ln\left(\frac{1}{\cos\theta}+\tan\theta\right)+C
\]

\noindent 这就是标准结论：$\displaystyle\int\frac{d\theta}{\cos\theta} = \ln\left|\frac{1}{\cos\theta}+\tan\theta\right|+C$。

\bigskip
\hrule
\bigskip

\ailine{你看，整个推导每一步都是你学过的东西——乘 $\cos\theta$、换元 $t=\sin\theta$、部分分式。每一步单独拿出来你都会，只是组合在一起是第一次见，想不到是完全正常的。这个结论记住就好，考试直接用。}

\ailine{现在最后一步，把 $\theta$ 换回 $x$——$\tan\theta = x/a$，那 $\dfrac{1}{\cos\theta}$ 等于什么？你刚才其实已经算过了。}

\userline{刚才算的 $\dfrac{a}{\cos\theta} = \sqrt{x^2+a^2}$，所以 $\dfrac{1}{\cos\theta} = \dfrac{\sqrt{x^2+a^2}}{a}$，代进去就是 $\ln\dfrac{x+\sqrt{x^2+a^2}}{a}+C$，常数吸收掉就是 $\ln(x+\sqrt{x^2+a^2})+C$}

\ailine{你已经完全掌握了！你看，整道题从头到尾就一个核心操作：看到 $x^2+a^2$ $\to$ 联想到 $1+\tan^2\theta = \frac{1}{\cos^2\theta}$ $\to$ 令 $x=a\tan\theta$ 消根号。这样你就拥有解决这一类根号下"平方加平方"积分的能力了，以后碰到直接 $\tan$ 换元，稳稳的。}

\bigskip
\hrule
\bigskip

\noindent{\large\textbf{情绪价值型风格核心原则}}

\begin{enumerate}
  \item \textbf{不评价难度}：不说"这题简单/善良"（暗示不会的人蠢），不说"这题很难"（让人觉得做不出来）。
  \item \textbf{归因于经验而非能力}：用"想不到是因为第一次见""这个直觉是做了足够多题之后才有的"，而非"你应该想到"。
  \item \textbf{鼓励落在当下}："你现在会了""你已经掌握了""这样你就拥有了解决这一类问题的能力"。
  \item \textbf{不跳步}：用户不会的中间步骤不丢一个公式让他记，而是完整推一遍。
  \item \textbf{肯定每一个正确的小步骤}："有这个思路很好""根号消掉了，这就是最关键的一步，你自己做出来了"。
  \item \textbf{错误时温和纠正}："这个方向是对的，不过……"而非"不对"。
\end{enumerate}

\end{document}
